% Л11 -- Безопасность: аутентификация и авторизация. Slides.
%
% Companion to the spoken script in ../L11-security.md, which is the source of
% truth: the slides carry diagrams, code, real terminal output and short
% marks, and the delivery lives in the script only -- this deck carries no
% speaker notes. Slide text is in Russian because that is the content;
% everything else in this file is English.
%
% Build:  make          (latexmk -pdf; engine: pdflatex)
%
% Every C# listing was compiled on dotnet SDK 10.0 in a scratch
% `dotnet new web -n HelloWorld`; terminal output on the slides is captured
% from runs of that project with tokens from `dotnet user-jwts`. See
% ../CLAUDE.md for the verification table.
%
% House rules: ../CLAUDE.md and ../../../CLAUDE.md. Theme: ../../../../theme/.

\documentclass[aspectratio=169,10pt,t]{beamer}

\usetheme{dotnet}
\usetikzlibrary{arrows.meta,positioning,calc,decorations.pathreplacing,fit}
\usepackage{array}   % ragged-right table columns

% A ragged-right paragraph column: justified cells in narrow columns hyphenate
% badly, and a split word is the one thing the back row cannot read.
\newcolumntype{L}[1]{>{\raggedright\arraybackslash}p{#1}}

\dotnetlecture{Л11}

\title[Л11. Безопасность: аутентификация и авторизация]{Безопасность:
  аутентификация и авторизация}
\author{Дмитрий Курочкин}
\date{}

% Code sizes, one step down from \footnotesize. Smaller than 6.6pt does not
% carry to the back of the room.
\newcommand{\codesm}{\ttfamily\fontsize{7.6}{9.0}\selectfont}
\newcommand{\codexs}{\ttfamily\fontsize{7.0}{8.3}\selectfont}
\newcommand{\codexxs}{\ttfamily\fontsize{6.8}{7.8}\selectfont}
\newcommand{\termsm}{\ttfamily\fontsize{6.6}{7.9}\selectfont}

% A mark at the start of a listing line, for labels set beside the listing.
% Used through the listing's escapechar; the labels are drawn by an overlay
% picture after the listing, so the build needs two passes (latexmk does that).
\newcommand{\mk}[1]{\tikz[remember picture,overlay]\coordinate (#1) at (0,0.55ex);}

\tikzset{
  rawbrace/.style={decorate,decoration={brace,amplitude=4pt},hairline,
                   line width=0.8pt},
  rawlab/.style={anchor=west,font=\footnotesize,inner sep=0pt},
  box/.style={draw=hairline,line width=0.8pt,fill=codebg,inner sep=4pt,
              align=center,minimum height=0.9cm},
  arr/.style={-{Straight Barb[length=1.6mm]},muted,line width=0.7pt},
  arra/.style={-{Straight Barb[length=1.8mm]},accent,line width=0.9pt},
  arrl/.style={-{Straight Barb[length=1.8mm]},alt,line width=0.9pt},
}

\begin{document}
\shorthandoff{"}

% ------------------------------------------------------------------ title --
% 1
\begin{frame}[plain]
  \titlepagednc{Лекция 11}{Безопасность: аутентификация и авторизация}{}{Дмитрий Курочкин}
\end{frame}

% ===================================== 1. AUTHENTICATION VS AUTHORIZATION ==
\section{1. Аутентификация vs авторизация}

% 2
\begin{frame}[plain]
  \sectionopener{1}{Аутентификация vs авторизация}{Сначала \hi{кто}, потом ---
    что можно.}
\end{frame}

% 3
\begin{frame}[plain]
  \vote{Аутентификация и авторизация. В чём разница?}
\end{frame}

% 4
\begin{frame}{Два вопроса}
  \slidetop
  \versus{Аутентификация}{%
    \begin{cbul}{accent}
      \item кто вы
      \item пароль, токен, сертификат
      \item проверяется доказательство
      \item отказ --- \hi{401}: «предъявите»
    \end{cbul}}{Авторизация}{%
    \begin{cbul}{alt}
      \item что вам можно
      \item правила для известной личности
      \item проверяется право
      \item отказ --- \lo{403}: предъявлять нечего
    \end{cbul}}

  \slidegap
  \hrulethin{0.42}
  \rulegap
  {\large Порядок жёсткий: нельзя решить, можно ли, не зная \hi{кому}.}
\end{frame}

% 5
\begin{frame}[fragile]{401 и 403}
  \slidetop
  \claim{Два разных ответа на два разных вопроса.}

\begin{lstlisting}[style=term,basicstyle=\termsm,escapechar=|]
$ curl -s -i http://localhost:5090/orders/1
|\mk{u1}|HTTP/1.1 401 Unauthorized
Content-Length: 0
Date: Thu, 12 Nov 2026 06:02:59 GMT
Server: Kestrel
WWW-Authenticate: Bearer

$ curl -s -i -H "Authorization: Bearer $TOKEN" http://localhost:5090/orders/7
|\mk{u2}|HTTP/1.1 403 Forbidden
Content-Length: 0
Date: Thu, 12 Nov 2026 06:02:59 GMT
Server: Kestrel
\end{lstlisting}
  \begin{tikzpicture}[remember picture,overlay,every node/.style={rawlab}]
    \node[text=accent] at ([xshift=5.0cm]u1) {не знаем, кто ты};
    \node[text=alt] at ([xshift=5.0cm]u2) {знаем --- и нельзя};
  \end{tikzpicture}

  \aside{401 на «нельзя» заставляет входить заново. 403 на «не знаем» не говорит,
    что нужно войти. Оба --- баги контракта.}
\end{frame}

% 6
\begin{frame}{Три слова фреймворка}
  \slidetop
  {\renewcommand{\arraystretch}{1.6}
  \begin{tabular}{@{}L{3.0cm}L{4.6cm}L{5.2cm}@{}}
    \hi{\textbf{схема}}     & способ доказать: имя и обработчик & токен в заголовке,
      cookie, сертификат \\
    \hi{\textbf{challenge}} & личности нет, а её требуют: «иди докажи» & токен --- 401;
      cookie --- страница входа \\
    \lo{\textbf{forbid}}    & личность есть, правило говорит «нельзя» & токен --- 403;
      cookie --- «доступ запрещён» \\
  \end{tabular}}

  \slidegap
  \aside{Код не пишет статус руками: он говорит «нужен пользователь», а форму
    ответа выбирает схема.}
\end{frame}

% 7
\begin{frame}{Где это в конвейере}
  \slidetop
  \claim{Авторизация работает с \hi{результатом} аутентификации.}

  \begin{center}
    \begin{tikzpicture}[font=\footnotesize,
      tb/.style={box,minimum width=3.3cm,minimum height=1.15cm}]
      \node[font=\footnotesize,text=muted] (rq) at (-2.6,0) {запрос};
      \node[tb] (a1) at (0,0)   {\textbf{аутентификация}\\схема заполняет \texttt{User}};
      \node[tb] (a2) at (4.3,0) {\textbf{авторизация}\\требования маршрута};
      \node[tb,draw=accent,line width=1.0pt] (a3) at (8.6,0)
        {\textbf{обработчик}\\только если всё прошло};
      \draw[arr] (rq) -- (a1);
      \draw[arr] (a1) -- (a2);
      \draw[arra] (a2) -- (a3);
      \node[font=\footnotesize,text=accent] (c1) at (2.9,-1.55) {challenge: 401};
      \node[font=\footnotesize,text=alt] (c2) at (5.7,-1.55) {forbid: 403};
      \draw[arra] ([xshift=-0.7cm]a2.south) -- (c1.north);
      \draw[arrl] ([xshift=0.7cm]a2.south) -- (c2.north);
    \end{tikzpicture}
  \end{center}

  \slidegap
  \aside{Доказательства нет --- \texttt{User} пустой, и это не ошибка: анонимный
    запрос нормален, пока никто не потребовал личность.}
\end{frame}

% 8
\begin{frame}[fragile]{Закрытость по умолчанию}
  \slidetop
\begin{lstlisting}[style=csharp,basicstyle=\codexs,escapechar=|]
builder.Services.AddAuthorizationBuilder()
|\mk{f1}|    .SetFallbackPolicy(new AuthorizationPolicyBuilder()
        .RequireAuthenticatedUser()
        .Build());

|\mk{f2}|app.MapGet("/health", () => "Healthy").AllowAnonymous();
|\mk{f3}|app.MapGet("/forgotten", () => "new route, nobody protected it");
\end{lstlisting}
  \begin{tikzpicture}[remember picture,overlay,every node/.style={rawlab}]
    \node[text=accent] at ([xshift=10.6cm]f1) {где не сказано иное};
    \node at ([xshift=10.6cm]f2) {открыт явно};
    \node at ([xshift=10.6cm]f3) {забыли защитить};
  \end{tikzpicture}

\begin{lstlisting}[style=term,basicstyle=\termsm]
$ curl -s -i http://localhost:5090/forgotten
HTTP/1.1 401 Unauthorized
Content-Length: 0
Date: Thu, 12 Nov 2026 06:03:00 GMT
Server: Kestrel
WWW-Authenticate: Bearer
\end{lstlisting}

  \aside{«Забыл открыть» видно сразу: 401. «Забыл закрыть» \lo{невидимо} до
    инцидента.}
\end{frame}

% ============================================================ 2. THE TOKEN ==
\section{2. Токен как способ остаться stateless}

% 9
\begin{frame}[plain]
  \sectionopener{2}{Токен как способ остаться stateless}{Токен --- подписанное
    \hi{утверждение}, которое клиент носит с собой.}
\end{frame}

% 10
\begin{frame}[plain]
  \vote{Три куска через точку: буквы, цифры, минусы. Это зашифрованные данные
    или закодированные?}
\end{frame}

% 11
\begin{frame}[fragile]{JWT: три части через точку}
  \slidetop
\begin{lstlisting}[style=term,basicstyle=\termsm,escapechar=`]
$ echo $TOKEN | tr . '\n' | fold -w 88
`\mk{w1}`eyJhbGciOiJIUzI1NiIsInR5cCI6IkpXVCJ9
eyJ1bmlxdWVfbmFtZSI6IjQyIiwic3ViIjoiNDIiLCJqdGkiOiJhNzUzOTkxMSIsInJvbGUiOiJjdXN0b21lciIs
ImF1ZCI6Imh0dHA6Ly9sb2NhbGhvc3Q6NTA5MCIsIm5iZiI6MTc5NDQ2MzM3MywiZXhwIjoxODAyNDEyMTczLCJp
`\mk{w2}`YXQiOjE3OTQ0NjMzNzMsImlzcyI6ImRvdG5ldC11c2VyLWp3dHMifQ
`\mk{w3}`O23hD9QOUmzjsrDBKtYY-zHTP2W0lyM_36qA19Ly4HA
\end{lstlisting}
  \begin{tikzpicture}[remember picture,overlay,every node/.style={rawlab}]
    \node[text=accent] at ([xshift=9.4cm]w1) {заголовок};
    \node[text=accent] at ([xshift=9.4cm]w2) {нагрузка};
    \node[text=alt] at ([xshift=9.4cm]w3) {подпись};
  \end{tikzpicture}

  \slidegap
  \begin{cbul}{accent}
    \item заголовок и нагрузка --- \hi{JSON}, закодированный в текст: буквы, цифры,
      минус и подчёркивание
    \item подпись --- байты: результат криптографической операции над первыми двумя
    \item около килобайта --- в \hi{каждом} запросе, в заголовке
  \end{cbul}
\end{frame}

% 12
\begin{frame}[fragile]{Закодирован --- не зашифрован}
  \vspace{0.3em}
\begin{lstlisting}[style=term,basicstyle=\termsm,escapechar=`]
$ echo $TOKEN | jq -R 'split(".")[0:2][] | @base64d | fromjson'
{
`\mk{h1}`  "alg": "HS256",
  "typ": "JWT"
}
{
  "unique_name": "42",
`\mk{h2}`  "sub": "42",
  "jti": "a7539911",
`\mk{h3}`  "role": "customer",
`\mk{h4}`  "aud": "http://localhost:5090",
  "nbf": 1794463373,
`\mk{h5}`  "exp": 1802412173,
  "iat": 1794463373,
`\mk{h6}`  "iss": "dotnet-user-jwts"
}
\end{lstlisting}
  \begin{tikzpicture}[remember picture,overlay,every node/.style={rawlab}]
    \node at ([xshift=7.2cm]h1) {алгоритм подписи};
    \node at ([xshift=7.2cm]h2) {субъект: о ком токен};
    \node at ([xshift=7.2cm]h3) {добавил издатель};
    \node at ([xshift=7.2cm]h4) {аудитория: для кого};
    \node at ([xshift=7.2cm]h5) {срок, секунды эпохи};
    \node at ([xshift=7.2cm]h6) {издатель: кто выдал};
  \end{tikzpicture}

  \aside{Раскодирует \lo{любой}, у кого есть строка: алфавит стандартный.}
\end{frame}

% 13
\begin{frame}{Утверждения}
  \slidetop
  \claim{Нагрузка --- пары имя--значение: \hi{claims}.}

  {\renewcommand{\arraystretch}{1.25}\small
  \begin{tabular}{@{}L{1.6cm}L{4.4cm}@{\hspace{0.9cm}}L{1.6cm}L{4.6cm}@{}}
    \hi{\texttt{iss}} & издатель: кто выдал          & \texttt{iat}   & момент выдачи \\
    \hi{\texttt{sub}} & субъект: о ком               & \texttt{nbf}   & не раньше \\
    \hi{\texttt{aud}} & аудитория: для кого          & \texttt{jti}   & идентификатор токена \\
    \hi{\texttt{exp}} & срок: до какого момента      & \texttt{scope} & области: для чего дан доступ \\
  \end{tabular}}

  \slidegap
  \aside{Остальное добавляет издатель: роли, отдел. Утверждение сделано \hi{в
    момент выдачи} --- и до обновления не меняется.}
\end{frame}

% 14
\begin{frame}[fragile]{Имена утверждений в .NET}
  \slidetop
  \claim{Короткие имена из токена становятся \lo{длинными}.}

\begin{lstlisting}[style=term,basicstyle=\termsm]
$ curl -s -H "Authorization: Bearer $TOKEN" http://localhost:5090/me | jq -c '.[]'
{"type":"http://schemas.xmlsoap.org/ws/2005/05/identity/claims/name","value":"42"}
{"type":"http://schemas.xmlsoap.org/ws/2005/05/identity/claims/nameidentifier","value":"42"}
{"type":"jti","value":"a7539911"}
{"type":"http://schemas.microsoft.com/ws/2008/06/identity/claims/role","value":"customer"}
{"type":"aud","value":"http://localhost:5090"}
{"type":"nbf","value":"1794463373"}
{"type":"exp","value":"1802412173"}
{"type":"iat","value":"1794463373"}
{"type":"iss","value":"dotnet-user-jwts"}
\end{lstlisting}

  \rulegap
  \aside{Ищут по константе --- \texttt{ClaimTypes.NameIdentifier}, а не по строке
    \texttt{sub}. Переименование отключается одной настройкой; в команде это
    решают \hi{один раз}.}
\end{frame}

% 15
\begin{frame}{Токен --- не профиль}
  \slidetop
  \versus{В токене}{%
    \begin{cbul}{accent}
      \item идентификатор
      \item роли, разрешения
      \item срок
      \item издатель, аудитория
    \end{cbul}}{У издателя и в базе}{%
    \begin{cbul}{alt}
      \item имя, почта, телефон
      \item отдел, должность
      \item лимит на заказы
      \item всё, что устаревает
    \end{cbul}}

  \slidegap
  \hrulethin{0.42}
  \rulegap
  \aside{Утверждений ровно столько, сколько нужно сервису для \hi{решений}.
    Чем больше в токене, тем больше того, что может устареть.}
\end{frame}

% 16
\begin{frame}[plain]
  \vote{Можно ли положить в токен номер карты, чтобы сервис оплаты не ходил за
    ним в базу?}
\end{frame}

% 17
\begin{frame}[plain]
  \statement{Токен \hi{подписывается}, а не шифруется. Подпись защищает от
    подделки, \lo{не от чтения}.}
\end{frame}

% 18
\begin{frame}[fragile]{Подпись}
  \slidetop
  \claim{Изменили символ --- подпись \lo{не сошлась}.}

\begin{lstlisting}[style=term,basicstyle=\termsm]
$ curl -s -i -H "Authorization: Bearer ${TOKEN/NDIi/NDMi}" http://localhost:5090/orders/1
HTTP/1.1 401 Unauthorized
Content-Length: 0
Date: Thu, 12 Nov 2026 06:03:00 GMT
Server: Kestrel
WWW-Authenticate: Bearer error="invalid_token", error_description="The signature key was not found"
\end{lstlisting}

  \rulegap
  \begin{cbul}{accent}
    \item \texttt{NDIi} $\to$ \texttt{NDMi} в нагрузке: \texttt{sub} был 42, стал 43
    \item проверяющий пересчитал подпись --- ни с одним его ключом она не сошлась
    \item сошлась --- нагрузка \hi{ровно такая}, какой её подписал издатель
  \end{cbul}
\end{frame}

% 19
\begin{frame}{Два вида ключей}
  \slidetop
  \versus{Асимметричный}{%
    \begin{cbul}{accent}
      \item пара ключей
      \item приватным подписывают
      \item публичным только проверяют
      \item приватный --- у издателя
      \item для нескольких сервисов
    \end{cbul}}{Симметричный}{%
    \begin{cbul}{alt}
      \item один секрет
      \item им подписывают
      \item им же и проверяют
      \item секрет --- у каждого
      \item для одного приложения
    \end{cbul}}

  \slidegap
  \hrulethin{0.42}
  \rulegap
  \aside{Алгоритм определяете \hi{вы}, а не токен: заголовок может написать кто
    угодно. Настройка --- список разрешённых алгоритмов.}
\end{frame}

% 20
\begin{frame}[fragile]{Откуда публичный ключ}
  \slidetop
\begin{lstlisting}[style=json,basicstyle=\codesm]
{ "alg": "RS256", "typ": "JWT", "kid": "2026-09-key-1" }
\end{lstlisting}

  \begin{center}
    \begin{tikzpicture}[font=\footnotesize,
      tb/.style={box,minimum width=4.0cm,minimum height=1.1cm}]
      \node[tb] (i) at (0,0) {\textbf{издатель}\\набор ключей (JWKS)};
      \node[tb,draw=accent,line width=1.0pt] (s) at (8.2,0)
        {\textbf{сервис}\\кэш публичных ключей};
      \draw[arra] (i) -- node[above,font=\scriptsize,text=accent] {ключ с таким
        \texttt{kid}} node[below,font=\scriptsize,text=muted] {один раз, потом из кэша} (s);
    \end{tikzpicture}
  \end{center}

  \rulegap
  \begin{cbul}{accent}
    \item незнакомый \texttt{kid} --- сервис идёт за набором заново
    \item издатель меняет ключи, \hi{не предупреждая} никого: ротация без простоя
    \item издатель недоступен --- известные ключи работают, незнакомый \texttt{kid}
      \lo{отклоняется}
  \end{cbul}
\end{frame}

% 21
\begin{frame}{Четыре проверки}
  \slidetop
  \begin{center}
    \begin{tikzpicture}[font=\footnotesize,
      tb/.style={box,minimum width=2.9cm,minimum height=1.25cm}]
      \node[tb] (t1) at (0,0)    {\textbf{подпись}\\не подделан};
      \node[tb] (t2) at (3.35,0) {\textbf{издатель}\\выдал тот,\\кому верим};
      \node[tb] (t3) at (6.7,0)  {\textbf{аудитория}\\выдан для нас};
      \node[tb] (t4) at (10.05,0) {\textbf{срок}\\ещё действует};
    \end{tikzpicture}
  \end{center}

  \slidegap
  {\large Все четыре, каждый запрос, \hi{без похода в базу}.}

  \slidegap
  \aside{Ключ в кэше, остальное --- математика и сравнение строк. На расхождение
    часов между машинами --- допуск в несколько минут.}
\end{frame}

% 22
\begin{frame}[fragile]{Отказ с причиной}
  \slidetop
\begin{lstlisting}[style=term,basicstyle=\termsm,escapechar=|]
$ curl -s -i -H "Authorization: Bearer $EXPIRED" http://localhost:5090/orders/1
HTTP/1.1 401 Unauthorized
Content-Length: 0
Date: Thu, 12 Nov 2026 06:03:00 GMT
Server: Kestrel
WWW-Authenticate: Bearer error="invalid_token",
error_description="The token expired at '11/11/2026 10:10:00'"

$ curl -s -i -H "Authorization: Bearer $REPORTS" http://localhost:5090/orders/1
HTTP/1.1 401 Unauthorized
Content-Length: 0
Date: Thu, 12 Nov 2026 06:03:00 GMT
Server: Kestrel
WWW-Authenticate: Bearer error="invalid_token",
error_description="The audience '(null)' is invalid"
\end{lstlisting}

  \rulegap
  \aside{Первый токен истёк вчера. Второй подписан \hi{тем же} издателем, но выдан
    сервису отчётов: чужая аудитория --- \lo{отказ}.}
\end{frame}

% 23
\begin{frame}{Stateless}
  \slidetop
  \claim{Любой экземпляр узнаёт, кто пришёл, \hi{никого не спрашивая}.}

  \begin{center}
    \begin{tikzpicture}[font=\footnotesize,
      tb/.style={box,minimum width=3.4cm,minimum height=0.72cm}]
      \node[font=\footnotesize,text=muted,align=center] (c) at (-0.6,0) {запрос\\с токеном};
      \node[tb,minimum width=2.8cm,minimum height=1.0cm] (lb) at (2.8,0) {балансировщик};
      \node[tb] (e1) at (7.2,0.95)  {экземпляр 1};
      \node[tb] (e2) at (7.2,0)     {экземпляр 2};
      \node[tb] (e3) at (7.2,-0.95) {экземпляр 3};
      \draw[arr] (c) -- (lb);
      \draw[arr] (lb.east) -- (e1.west);
      \draw[arr] (lb.east) -- (e2.west);
      \draw[arr] (lb.east) -- (e3.west);
      \draw[rawbrace] ([xshift=0.25cm]e1.north east) -- ([xshift=0.25cm]e3.south east);
      \node[rawlab,text=accent] at ([xshift=0.6cm]e2.east) {тот же публичный ключ};
    \end{tikzpicture}
  \end{center}

  \slidegap
  \aside{Ни общего хранилища сессий, ни липких сессий. Десятый экземпляр сразу
    принимает те же токены.}
\end{frame}

% 24
\begin{frame}[plain]
  \vote{Пользователя уволили в три часа дня. Его токен выдан в два, срок --- до
    пяти. Что с его запросами с трёх до пяти?}
\end{frame}

% 25
\begin{frame}{Мгновенного отзыва нет}
  \slidetop
  \versus{JWT}{%
    \begin{cbul}{accent}
      \item содержимое внутри
      \item проверка без издателя
      \item отозвать нельзя: ждём срока
      \item сервис независим
    \end{cbul}}{Непрозрачный токен}{%
    \begin{cbul}{alt}
      \item случайная строка
      \item каждый запрос --- к издателю
      \item отзыв мгновенный
      \item издатель --- узкое место
    \end{cbul}}

  \slidegap
  \hrulethin{0.42}
  \rulegap
  \aside{Отозвать --- значит где-то \lo{помнить}. Список отозванных на каждом
    запросе --- то же хранилище, только с лишним слоем.}
\end{frame}

% 26
\begin{frame}{Два токена}
  \slidetop
  \versus{Access}{%
    \begin{cbul}{accent}
      \item минуты: 5, 10, 15
      \item stateless
      \item предъявляют сервису
      \item доживает свой срок
    \end{cbul}}{Refresh}{%
    \begin{cbul}{alt}
      \item дни или недели
      \item хранится у издателя
      \item предъявляют только издателю
      \item отзывается одной записью
    \end{cbul}}

  \slidegap
  \hrulethin{0.42}
  \rulegap
  {\large Уволили $\to$ refresh отозван $\to$ через \hi{минуты} доступ закрыт.}
\end{frame}

% 27
\begin{frame}[plain]
  \vote{Часы на сервере отстают на десять минут. Срок токена --- пять минут. Что
    увидит сервис?}
\end{frame}

% 28
\begin{frame}[fragile]{Включить в сервисе}
  \vspace{0.3em}
\begin{lstlisting}[style=csharp,basicstyle=\codesm,escapechar=|]
builder.Services
|\mk{b1}|    .AddAuthentication(JwtBearerDefaults.AuthenticationScheme)
    .AddJwtBearer(options =>
    {
|\mk{b2}|        options.Authority =
            builder.Configuration["Auth:Authority"];
|\mk{b3}|        options.Audience =
            builder.Configuration["Auth:Audience"];
    });

builder.Services.AddAuthorization();

var app = builder.Build();

|\mk{b4}|var orders = app.MapGroup("/orders").RequireAuthorization();
\end{lstlisting}
  \begin{tikzpicture}[remember picture,overlay,every node/.style={rawlab}]
    \node at ([xshift=10.3cm]b1) {схема по умолчанию};
    \node[text=accent] at ([xshift=10.3cm]b2) {издатель};
    \node[text=accent] at ([xshift=10.3cm]b3) {аудитория --- мы};
    \node at ([xshift=10.3cm]b4) {любой, но известный};
  \end{tikzpicture}

\end{frame}

% 29
\begin{frame}[fragile]{Пока издателя нет}
  \vspace{0.3em}
\begin{lstlisting}[style=term,basicstyle=\termsm]
$ dotnet user-jwts create --name 42 --role customer
New JWT saved with ID 'a7539911'.
Name: 42
Roles: [customer]

Token: eyJhbGciOiJIUzI1NiIsInR5cCI6IkpXVCJ9.eyJ1bmlxdWVfbmFtZSI6IjQyIiwic3ViIjoiNDIiLCJqdGki...
$ curl -s -i -H "Authorization: Bearer $TOKEN" http://localhost:5090/orders/1
HTTP/1.1 200 OK
Content-Type: application/json; charset=utf-8
Date: Thu, 12 Nov 2026 06:03:00 GMT
Server: Kestrel
Transfer-Encoding: chunked

{"id":1,"status":"Paid","total":25.00}
\end{lstlisting}

  \rulegap
  \aside{\texttt{user-jwts}: ключ подписи --- в секретах пользователя, издатель и
    аудитория --- в конфигурации разработки. Авторизацию отлаживают \hi{до}
    настоящего издателя.}
\end{frame}

% 30
\begin{frame}{Что за двумя строками}
  \slidetop
  \claim{«Почти без кода» --- не «почти ничего не происходит».}

  {\renewcommand{\arraystretch}{1.6}
  \begin{tabular}{@{}L{3.2cm}L{9.6cm}@{}}
    \hi{\textbf{один раз}}      & по адресу издателя --- документ с описанием; по ссылке
      из него --- набор ключей; в кэш \\
    \hi{\textbf{каждый запрос}} & токен из заголовка $\to$ ключ по \texttt{kid} $\to$
      подпись, издатель, аудитория, срок \\
    \hi{\textbf{удача}}         & утверждения становятся \texttt{User} в контексте \\
    \lo{\textbf{неудача}}       & \texttt{User} пустой; группа потребует его --- challenge, 401 \\
  \end{tabular}}

  \slidegap
  \aside{В обработчик пользователь приходит параметром по типу:
    \texttt{ClaimsPrincipal}.}
\end{frame}

% 31
\begin{frame}[plain]
  \keyline{Токен --- подписанные утверждения: база не нужна.\\
    Читается всеми --- в нём нет чувствительного.\\
    Не отзывается --- поэтому короткий.}
\end{frame}

% ========================================================= 3. ACCESS RULES ==
\section{3. Правила доступа}

% 32
\begin{frame}[plain]
  \sectionopener{3}{Правила доступа}{Доступ к \hi{ресурсу}, а не только к
    маршруту.}
\end{frame}

% 33
\begin{frame}{Три ступени}
  \slidetop
  {\renewcommand{\arraystretch}{1.5}
  \begin{tabular}{@{}L{3.4cm}L{4.7cm}L{4.7cm}@{}}
    & \colhead{accent}{что это} & \colhead{alt}{когда} \\
    \textbf{роли}          & слово в токене, одна строка на маршруте & правил не больше,
      чем ролей \\
    \textbf{политики}      & именованное правило из условий & правило меняется, маршруты
      не тронуты \\
    \hi{\textbf{требования и обработчики}} & данные и код, зависимости из контейнера
      & условие требует кода: возраст, отдел, владелец \\
  \end{tabular}}

  \slidegap
  \aside{Минимальные права: украденный токен с минимальными правами ---
    \hi{минимальный} ущерб.}
\end{frame}

% 34
\begin{frame}[fragile]{Политика}
  \slidetop
\begin{lstlisting}[style=csharp,basicstyle=\codesm,escapechar=|]
builder.Services.AddAuthorizationBuilder()
|\mk{p1}|    .AddPolicy("CanManageOrders", policy => policy
|\mk{p2}|        .RequireRole("manager")
|\mk{p3}|        .RequireClaim("department", "sales", "support"));

orders.MapGet("/", (IOrderService service) =>
        service.All().Select(o => o.ToDto()))
|\mk{p4}|    .RequireAuthorization("CanManageOrders");
\end{lstlisting}
  \begin{tikzpicture}[remember picture,overlay,every node/.style={rawlab}]
    \node[text=accent] at ([xshift=9.9cm]p1) {имя --- один раз};
    \draw[rawbrace] ([xshift=9.65cm,yshift=0.35em]p2) -- ([xshift=9.65cm,yshift=-0.3em]p3);
    \node at ([xshift=9.9cm]$(p2)!0.5!(p3)$) {условия --- как «и»};
    \node at ([xshift=9.9cm]p4) {маршрут знает имя};
  \end{tikzpicture}

\begin{lstlisting}[style=term,basicstyle=\termsm]
$ curl -s -w '\n%{http_code}\n' -H "Authorization: Bearer $TOKEN" http://localhost:5090/orders

403
$ curl -s -w '\n%{http_code}\n' -H "Authorization: Bearer $MANAGER" http://localhost:5090/orders
[{"id":1,"status":"Paid","total":25.00},{"id":7,"status":"New","total":12.50}]
200
\end{lstlisting}
\end{frame}

% 35
\begin{frame}[plain]
  \vote{Маршрут защищён. Пользователь 42 запрашивает заказ 7 --- заказ
    пользователя 9. Что вернёт сервис?}
\end{frame}

% 36
\begin{frame}{Две проверки}
  \slidetop
  \versus{На маршруте}{%
    \begin{cbul}{accent}
      \item кто может вызывать \texttt{GET /orders/\{id\}}
      \item видит пользователя и маршрут
      \item до обработчика: заказа ещё нет
      \item ловит анонимного
    \end{cbul}}{На ресурсе}{%
    \begin{cbul}{alt}
      \item кому можно заказ 7
      \item видит пользователя и заказ
      \item после загрузки
      \item ловит чужого
    \end{cbul}}

  \slidegap
  \hrulethin{0.42}
  \rulegap
  {\large Только первая --- \lo{уязвимость}: поменял цифру в адресе --- получил
    чужой заказ.}
\end{frame}

% 37
\begin{frame}[fragile]{Требование и обработчик}
  \vspace{0.3em}
\begin{lstlisting}[style=csharp,basicstyle=\codexs,escapechar=`]
`\mk{r1}`public sealed class OrderOwnerRequirement
    : IAuthorizationRequirement;

public sealed class OrderOwnerHandler
`\mk{r2}`    : AuthorizationHandler<OrderOwnerRequirement, Order>
{
    protected override Task HandleRequirementAsync(
        AuthorizationHandlerContext context,
        OrderOwnerRequirement requirement, Order order)
    {
        var userId = context.User.FindFirstValue(ClaimTypes.NameIdentifier);
`\mk{r3}`        if (userId == order.CustomerId
            || context.User.IsInRole("admin"))
`\mk{r4}`            context.Succeed(requirement);
        return Task.CompletedTask;
    }
}
\end{lstlisting}
  \begin{tikzpicture}[remember picture,overlay,every node/.style={rawlab}]
    \node at ([xshift=10.2cm]r1) {требование --- данные};
    \node[text=accent] at ([xshift=10.2cm]r2) {и тип ресурса};
    \node at ([xshift=10.2cm]r3) {правило --- кодом};
    \node[text=alt] at ([xshift=10.2cm]r4) {промолчал --- отказ};
  \end{tikzpicture}
\end{frame}

% 38
\begin{frame}[fragile]{Проверка на ресурсе}
  \vspace{0.3em}
\begin{lstlisting}[style=csharp,basicstyle=\codexs,escapechar=`]
orders.MapGet("/{id:int}",
`\mk{e1}`    async Task<Results<Ok<OrderDto>, NotFound, ForbidHttpResult>> (
        int id, ClaimsPrincipal user, IOrderService service,
        IAuthorizationService auth, CancellationToken ct) =>
{
`\mk{e2}`    var order = await service.FindAsync(id, ct);
    if (order is null) return TypedResults.NotFound();
`\mk{e3}`    var result = await auth.AuthorizeAsync(
        user, order, new OrderOwnerRequirement());
    return result.Succeeded
        ? TypedResults.Ok(order.ToDto())
`\mk{e4}`        : TypedResults.Forbid();
});
\end{lstlisting}
  \begin{tikzpicture}[remember picture,overlay,every node/.style={rawlab}]
    \node at ([xshift=10.2cm]e1) {три исхода --- в типе};
    \node at ([xshift=10.2cm]e2) {сначала загрузили};
    \node[text=accent] at ([xshift=10.2cm]e3) {потом проверили};
    \node[text=alt] at ([xshift=10.2cm]e4) {forbid: 403};
  \end{tikzpicture}

  \slidegap
  \aside{Проверка на ресурсе \hi{не отменяет} проверку на маршруте: анонимный
    запрос до загрузки заказа не доходит.}
\end{frame}

% 39
\begin{frame}{403 или 404 для чужого}
  \slidetop
  \versus{403}{%
    \begin{cbul}{accent}
      \item строго по смыслу: есть, но нельзя
      \item подтверждает: номер занят
      \item перебором видно, сколько заказов
    \end{cbul}}{404}{%
    \begin{cbul}{alt}
      \item «для тебя его нет»
      \item существование не раскрыто
      \item перебор ничего не даёт
    \end{cbul}}

  \slidegap
  \hrulethin{0.42}
  \rulegap
  \aside{Осознанный выбор --- \hi{один} на весь сервис, а не в каждом маршруте
    по-разному.}
\end{frame}

% 40
\begin{frame}[fragile]{Свой, чужой, администратор}
  \vspace{0.3em}
\begin{lstlisting}[style=term,basicstyle=\termsm,escapechar=|]
$ curl -s -w '\n%{http_code}\n' -H "Authorization: Bearer $TOKEN" http://localhost:5090/orders/1
{"id":1,"status":"Paid","total":25.00}
|\mk{o1}|200
$ curl -s -w '\n%{http_code}\n' -H "Authorization: Bearer $TOKEN" http://localhost:5090/orders/7

|\mk{o2}|403
$ curl -s -w '\n%{http_code}\n' -H "Authorization: Bearer $ADMIN" http://localhost:5090/orders/7
{"id":7,"status":"New","total":12.50}
|\mk{o3}|200
\end{lstlisting}
  \begin{tikzpicture}[remember picture,overlay,every node/.style={rawlab}]
    \node[text=accent] at ([xshift=1.6cm]o1) {свой};
    \node[text=alt] at ([xshift=1.6cm]o2) {чужой};
    \node at ([xshift=1.6cm]o3) {администратор};
  \end{tikzpicture}

  \rulegap
  \begin{cbul}{accent}
    \item один объект --- обработчик; список «мои заказы» --- \hi{фильтр} по владельцу
      в запросе к базе
    \item заказ загружают один раз: и для проверки, и для ответа
    \item правило владения --- в \hi{одном} обработчике, для всех маршрутов с заказом
  \end{cbul}
\end{frame}

% 41
\begin{frame}[plain]
  \keyline{Правило на маршруте --- кто может спрашивать.\\
    Правило на ресурсе --- кому можно этот объект.\\
    Только первое --- уязвимость.}
\end{frame}

% =================================================== 4. EXTERNAL PROVIDERS ==
\section{4. Внешние провайдеры}

% 42
\begin{frame}[plain]
  \sectionopener{4}{Внешние провайдеры}{Сервис не изобретает identity --- он к
    ней \hi{подключается}.}
\end{frame}

% 43
\begin{frame}[plain]
  \vote{Вы нажали «Войти через Google» на чужом сайте. Сайт узнал ваш пароль?}
\end{frame}

% 44
\begin{frame}{Четыре роли OAuth 2.0}
  \slidetop
  \begin{center}
    \begin{tikzpicture}[font=\footnotesize,
      tb/.style={box,minimum width=4.3cm,minimum height=1.15cm}]
      \node[tb] (ro) at (0,0)    {\textbf{владелец ресурса}\\пользователь};
      \node[tb] (as) at (8.4,0)  {\textbf{сервер авторизации}\\выдаёт токены};
      \node[tb] (cl) at (0,-2.3) {\textbf{клиент}\\приложение, которому нужен доступ};
      \node[tb,draw=accent,line width=1.0pt] (rs) at (8.4,-2.3)
        {\textbf{ресурсный сервер}\\наш сервис};
      \draw[arrl] (ro) -- node[above,font=\scriptsize,text=alt] {пароль --- только сюда} (as);
      \draw[arr] (as) -- node[above,sloped,font=\scriptsize,text=muted] {токены} (cl);
      \draw[arra] (cl) -- node[above,font=\scriptsize,text=accent] {токен доступа} (rs);
    \end{tikzpicture}
  \end{center}

  \slidegap
  \aside{Наш сервис не проверяет пароли и не знает пользователей: он принимает
    токены и \hi{проверяет} их.}
\end{frame}

% 45
\begin{frame}{Поток с кодом авторизации}
  \slidetop
  \claim{Пароль видел только \hi{сервер авторизации}.}

  {\renewcommand{\arraystretch}{1.3}
  \begin{tabular}{@{}L{0.7cm}L{12.1cm}@{}}
    \hi{\textbf{1}} & клиент отправляет пользователя на сервер авторизации \\
    \hi{\textbf{2}} & пользователь вводит пароль \hi{там} и соглашается дать доступ \\
    \hi{\textbf{3}} & сервер авторизации возвращает его к клиенту с одноразовым кодом \\
    \hi{\textbf{4}} & клиент обменивает код на токены --- напрямую \\
    \hi{\textbf{5}} & клиент ходит к нашему сервису с токеном доступа \\
  \end{tabular}}

  \slidegap
  \aside{Проверочный ключ: отпечаток случайной строки --- с первым запросом, сама
    строка --- с обменом. Перехваченный код \lo{не обменять}.}
\end{frame}

% 46
\begin{frame}{Потоки}
  \slidetop
  {\renewcommand{\arraystretch}{1.6}
  \begin{tabular}{@{}L{4.6cm}L{8.2cm}@{}}
    \hi{\textbf{код авторизации}}         & есть пользователь и браузер \\
    \hi{\textbf{учётные данные клиента}}  & пользователя нет: сервис вызывает сервис от
      своего имени \\
    \textbf{устройство}                   & браузера нет: код вводят на другом
      устройстве \\
    \textbf{обмен токенов}                & цепочка сервисов: дальше --- тоже от имени
      пользователя \\
  \end{tabular}}

  \slidegap
  \aside{Субъект в токене сервиса --- сервис, а не человек, и правила доступа для
    него \hi{свои}.}
\end{frame}

% 47
\begin{frame}{OpenID Connect: два токена}
  \slidetop
  \versus{Токен доступа}{%
    \begin{cbul}{accent}
      \item для ресурсного сервера
      \item с ним ходят к нам
      \item что клиенту можно
      \item аудитория --- наш сервис
    \end{cbul}}{Токен идентичности}{%
    \begin{cbul}{alt}
      \item для клиента
      \item к нам с ним нельзя
      \item кто вошёл: имя в углу экрана
      \item аудитория --- не наша
    \end{cbul}}

  \slidegap
  \hrulethin{0.42}
  \rulegap
  \aside{OAuth --- про доступ. OpenID Connect добавляет личность: токен
    идентичности, профиль и \hi{документ обнаружения} по известному адресу.}
\end{frame}

% 48
\begin{frame}{Подключиться к чему угодно}
  \slidetop
  \begin{center}
    \begin{tikzpicture}[font=\footnotesize,
      tb/.style={box,minimum width=3.5cm,minimum height=0.85cm}]
      \node[tb,draw=accent,line width=1.0pt,minimum width=3.3cm,minimum height=1.5cm]
        (s) at (0,0) {\textbf{наш сервис}\\\texttt{Authority}\\\texttt{Audience}};
      \node[tb] (k) at (-5.2,0.85)  {Keycloak};
      \node[tb] (i) at (-5.2,-0.85) {IdentityServer};
      \node[tb] (c) at (5.2,0.85)   {облачный провайдер};
      \node[tb] (t) at (5.2,-0.85)  {тестовый из SDK};
      \draw[arra] (k.east) -- (s);
      \draw[arra] (i.east) -- (s);
      \draw[arra] (c.west) -- (s);
      \draw[arra] (t.west) -- (s);
    \end{tikzpicture}
  \end{center}

  \slidegap
  {\large Меняется \hi{конфигурация}, не код.}

  \slidegap
  \aside{Один протокол, одна структура токена, один стандарт публикации ключей.
    Сервис, который сам хранит пароли, привязан к своей базе пользователей
    \lo{навсегда}.}
\end{frame}

% 49
\begin{frame}{Что настраиваем}
  \slidetop
  {\renewcommand{\arraystretch}{1.6}
  \begin{tabular}{@{}L{3.6cm}L{9.2cm}@{}}
    \hi{\textbf{издатель}}         & адрес сервера авторизации: документ обнаружения,
      имя издателя, набор ключей \\
    \hi{\textbf{аудитория}}        & имя сервиса у издателя: токен для соседнего
      сервиса отвергается \\
    \hi{\textbf{правила проверки}} & допуск на часы, алгоритмы подписи, откуда брать
      роли \\
  \end{tabular}}

  \slidegap
  \aside{Роли каждый издатель кладёт по-своему: их приводит к стандартным
    \hi{один класс} --- преобразование утверждений. Адрес издателя --- только HTTPS:
    подменили его --- подменили всех, кому сервис верит.}
\end{frame}

% 50
\begin{frame}[plain]
  \vote{А почему не просто ключ API --- строка в заголовке, которую сервис
    сверяет с базой?}
\end{frame}

% 51
\begin{frame}{Ключ API против токена}
  \slidetop
  \versus{Токен}{%
    \begin{cbul}{accent}
      \item живёт минуты, обновляется
      \item личность пользователя
      \item отзыв через издателя
      \item единый вход
    \end{cbul}}{Ключ API}{%
    \begin{cbul}{alt}
      \item выдан на месяцы
      \item только приложение
      \item утечка --- доступ на месяцы
      \item ни срока, ни обновления
    \end{cbul}}

  \slidegap
  \hrulethin{0.42}
  \rulegap
  \aside{Машина к машине, с ротацией и ограниченными правами --- ключ допустим.
    Там, где есть пользователь, --- шаг назад на \lo{пятнадцать лет}.}
\end{frame}

% 52
\begin{frame}[plain]
  \keyline{Пароли и личности --- чужая работа по стандарту.\\
    Сервис принимает токены.\\
    Издателя меняет конфигурация, а не код.}
\end{frame}

% ================================================== 5. SECRETS AND HYGIENE ==
\section{5. Секреты и гигиена}

% 53
\begin{frame}[plain]
  \sectionopener{5}{Секреты и гигиена}{Не в репозитории, не в логах, не в
    \hi{токене}.}
\end{frame}

% 54
\begin{frame}[fragile]{Локально --- секреты пользователя}
  \slidetop
\begin{lstlisting}[style=term,basicstyle=\termsm,escapechar=|]
|\mk{s1}|$ dotnet user-secrets init
Set UserSecretsId to 'ebb9ffa5-e2b3-43da-9de9-d8ef98bb5d31' for MSBuild project
'/home/user/HelloWorld/HelloWorld.csproj'.
|\mk{s2}|$ dotnet user-secrets set "Auth:ClientSecret" "dev-only-value"
Successfully saved Auth:ClientSecret to the secret store.
|\mk{s3}|$ dotnet user-secrets list
Auth:ClientSecret = dev-only-value
\end{lstlisting}
  \begin{tikzpicture}[remember picture,overlay,every node/.style={rawlab}]
    \node at ([xshift=9.6cm]s1) {идентификатор проекта};
    \node[text=accent] at ([xshift=9.6cm]s2) {ключ --- с двоеточиями};
    \node at ([xshift=9.6cm]s3) {что лежит};
  \end{tikzpicture}

  \slidegap
  \begin{cbul}{accent}
    \item значение --- в домашней папке, \hi{вне репозитория}; у коллеги своё
    \item читается только в разработке
    \item в файле настроек в репозитории --- ключ с пустым значением: видна структура
  \end{cbul}
\end{frame}

% 55
\begin{frame}{Где секретам место}
  \slidetop
  {\renewcommand{\arraystretch}{1.6}
  \begin{tabular}{@{}L{3.2cm}L{9.6cm}@{}}
    \hi{\textbf{локально}} & секреты пользователя \\
    \hi{\textbf{сборка}}   & хранилище секретов системы сборки: подставляет в
      переменные окружения шага, маскирует в выводе \\
    \hi{\textbf{прод}}     & переменные окружения или хранилище секретов: журнал
      чтений, ротация, ещё один провайдер конфигурации \\
  \end{tabular}}

  \slidegap
  \aside{Секрет, который не менялся два года, знают все, кто работал в компании за
    эти два года. Лучшие учётные данные --- \hi{короткоживущие}: истекают сами.}
\end{frame}

% 56
\begin{frame}[plain]
  \vote{Вы закоммитили строку подключения с паролем, через минуту удалили и
    закоммитили снова. Пароль в безопасности?}
\end{frame}

% 57
\begin{frame}{Три места, где секретов нет}
  \slidetop
  {\renewcommand{\arraystretch}{1.6}
  \begin{tabular}{@{}L{3.2cm}L{9.6cm}@{}}
    \lo{\textbf{репозиторий}} & он в истории. Попавший секрет считается утёкшим:
      его не удаляют, а \hi{меняют} \\
    \lo{\textbf{логи}}        & заголовок \texttt{Authorization}, тело запроса на
      вход, исключение со строкой подключения, токен в адресе \\
    \lo{\textbf{токен}}       & читается всеми: идентификатор вместо имени, роль
      вместо должности \\
  \end{tabular}}

  \slidegap
  \aside{Логировать можно и нужно: \hi{идентификатор} пользователя и идентификатор
    трассировки. Иначе не ответить, кто удалил заказ.}
\end{frame}

% 58
\begin{frame}[plain]
  \keyline{Секрет, попавший в репозиторий, в лог или в токен, --- уже не
    секрет.\\ Его место --- секреты пользователя и хранилище.}
\end{frame}

\end{document}
